\documentclass{beamer}
\usepackage{../tut-slides}
\usepackage{../mathoperatorsAuD}

\usepackage{csquotes}
\usepackage{cancel}

\usepackage{amsmath,amssymb}

\usepackage{tikz}
\usetikzlibrary{positioning,automata, matrix, trees}
\usetikzlibrary{calc,positioning,backgrounds,arrows.meta}
\usepackage{forest}


\usepackage{booktabs}
\usepackage{tabularx}
\usepackage{tabu}

\renewcommand{\tabularxcolumn}[1]{>{\hspace{0pt}}m{#1}}

\usepackage{listings}
\lstset{ 
	basicstyle=\footnotesize\ttfamily,        % the size of the fonts that are used for the code
	breakatwhitespace=false,         % sets if automatic breaks should only happen at whitespace
	breaklines=true,                 % sets automatic line breaking
	commentstyle=\itshape,    	     % comment style
	escapeinside={\%*}{*)},          % if you want to add LaTeX within your code
	extendedchars=true,              % lets you use non-ASCII characters; for 8-bits encodings only, does not work with UTF-8
	firstnumber=1,                % start line enumeration with line 1000
	frame=single,
	keywordstyle=\bfseries,       % keyword style
	morekeywords={}, 
	language=Prolog,                 % the language of the code
	numbers=left,                    % where to put the line-numbers; possible: (none, left, right)
	numbersep=5pt,                   % how far the line-numbers are from the code
	numberstyle=\tiny\color{cdgray!50}, % the style that is used for the line-numbers
	rulecolor=\color{cddarkblue}, 
	tabsize=2,	                   % sets default tabsize to 2 spaces
}
\lstdefinestyle{refutation}{ 
	basicstyle=\footnotesize\ttfamily,        % the size of the fonts that are used for the code
	breakatwhitespace=false,         % sets if automatic breaks should only happen at whitespace
	breaklines=true,                 % sets automatic line breaking
	commentstyle=\color{cdgray},    	     % comment style
	escapeinside={(*@}{@*)},          % if you want to add LaTeX within your code
	extendedchars=true,              % lets you use non-ASCII characters; for 8-bits encodings only, does not work with UTF-8
	firstnumber=1,                % start line enumeration with line 1000
	frame=none,
	keywordstyle=\bfseries,       % keyword style
	morekeywords={}, 
	language=Prolog,                 % the language of the code
	numbers=none,                    % where to put the line-numbers; possible: (none, left, right)
	numbersep=5pt,                   % how far the line-numbers are from the code
	numberstyle=\tiny\color{cdgray!50}, % the style that is used for the line-numbers
	rulecolor=\color{cddarkblue}, 
	tabsize=2,	                   % sets default tabsize to 2 spaces
}

\usepackage{textgreek}
\usepackage{adjustbox}


\renewcommand{\emph}[1]{\textbf{#1}}
\newcommand{\coloremph}[1]{\textcolor{cdpurple}{#1}}
\newcommand{\col}[1]{\textcolor{cdpurple}{\boldsymbol{#1}}}
\newcommand{\coll}[1]{\textcolor{cddarkgreen}{\boldsymbol{#1}}}
\newcommand{\colll}[1]{\textcolor{cdorange}{\boldsymbol{#1}}}
%\newcommand{\step}[2][]{\ensuremath{\overset{{#1} (\text{#2})}{=}}}
%\newcommand*{\astep}[2][]{\ensuremath{\overset{{#1} (\text{#2})}&{=}}}

\newcommand{\num}[1]{\ensuremath{\langle #1 \rangle}}


\pgfdeclareverticalshading{rainbow}{100bp}
{color(0bp)=(red); color(25bp)=(red); color(35bp)=(yellow);
	color(45bp)=(green); color(55bp)=(cyan); color(65bp)=(blue);
	color(75bp)=(violet); color(100bp)=(violet)}


\begin{document}	
	\title{Programmierung}
	\subtitle{Übung 9: Logikprogrammierung mit Prolog${}^-$}
	\author{Eric Kunze}
	\email{eric.kunze@tu-dresden.de}
	\city{TU Dresden}
	\date{13. Juni 2022}
%	\institute{Lehrstuhl für Grundlagen der Programmierung}
	\titlegraphic{\includegraphics[width=2cm]{../TUD-white.pdf}}
	
	\maketitle
	

%%%%%%%%%%%%%%%%%%%%%%%%%%%%%%%%%%%%%%%%%%%%%%%%%%%%%%%%%%%%%%%%%%%%%%%%%%%%%

{
\setbeamertemplate{frametitle}[pride]

\begin{frame} \frametitle{Pride Month}
	\centering
	\includegraphics[width=0.7\linewidth]{bmw-1.png}
	\pause
	\includegraphics[width=0.6\linewidth]{bmw-2.png}
	\begin{tikzpicture}[remember picture, overlay]
		\node [
		font=\tiny\color{cdgray!80},  % Einstellungen für die Schrift
		inner xsep=0,       % Abstand des Textes von unten
		% maximale Textbreite = Papierhöhe - 2*Abstand des Textes von unten:
		%			text width={\dimexpr\paperheight-2\footskip\relax},
		align=center,
		%			minimum height=8mm,% Breite des Randstreifens
		anchor=south west,
		] at ($(current page.south west)+(+10mm,+2.5mm)$)
		{Quelle: \url{https://twitter.com/BMW}};
	\end{tikzpicture}	
\end{frame}

\begin{frame} \frametitle{Beratungsmöglichkeiten} 
	\scriptsize

	\textbf{Queere Peerberatung:} Fay Uhlmann, \hyperref{mailto:queerpeer@tu-dresden.de}{}{}{queerpeer@tu-dresden.de}

	\begin{quote}
		Das Ziel der queeren Peerberatung ist, eine zentrale Anlaufstelle für queere Unimitglieder zu bieten. Sie ist ein offenes Ohr, um auf Augenhöhe über Diskriminierungserfahrungen aufgrund von Queerness, Problemen im Alltag, Outings, Transitionen, Sexualität, uvm. reden zu können. Auch soll sie bei der Bewältigung der vielen Hürden unterstützen, mit denen queere Personen noch immer im Alltag zu kämpfen haben.
	\end{quote}

	\textbf{Psychosoziale Beratungsstelle:} Studierendenwerk
	\begin{quote}
		[...]
		Bereiche des psychologischen Beratungsbedarfs können z.B. sein: [...]
		\begin{itemize}
			\item mangelndes Selbstwertgefühl
			\item Probleme im sozialen Umfeld
			\item Probleme mit Alkohol, Drogen, Online-Sucht
			\item depressive Verstimmungen
		\end{itemize}
		Bei Partnerschaftsproblemen ist auch Paarberatung möglich.
	\end{quote}
	{\tiny \url{https://www.studentenwerk-dresden.de/soziales/psychosoziale-beratung.html}}
\end{frame}

\begin{frame} \frametitle{Hilfe, Interaktion, \dots}
	\scriptsize
	
	\textbf{Gerede e.V.:} \url{https://gerede-dresden.de/}
	\begin{quote}
		Der Gerede e.V. versteht sich als Interessenvertretung für Schwule, Lesben, Bisexuelle, Transidente und Menschen mit vielfältigen Lebensweisen sowie deren Angehörige. 
		\begin{itemize}
			\item Beratungs-, Begegnungs- und Bildungsangebote sowie kulturelle Veranstaltungen
			\item für Jugendliche und junge Erwachsene, aber auch Eltern, Lehrer*innen, Erzieher*innen und Sozialpädagogen*innen
		\end{itemize}
	\end{quote}
	
		
	\textbf{Nightline Dresden:} \url{https://nightline-dresden.de/}
	\begin{quote}
		Egal ob Prüfungsangst, Zweifel an der Wahl des Studienfachs, Stress in der WG, Heimweh oder Liebeskummer – die Liste studentischer Sorgen ist lang und gerade bei solchen Problemen sind wir für Dich da. Auch wenn Du einfach mal reden willst, haben wir ein offenes Ohr für Dich.
		
		Wir sind während der Vorlesungszeit Dienstag, Donnerstag, Samstag und Sonntag von 21 bis 24 Uhr erreichbar!
		
		Telefonnummer: 0351 4277345		
	\end{quote}

	\textbf{Dresdner Wegweiser für Krisen- und Notsituationen:} \url{www.dresden.de/krisenwegweiser}
\end{frame}

}



\section{Logikprogrammierung und Prolog${}^-$}

%%%%%%%%%%%%%%%%%%%%%%%%%%%%%%%%%%%%%%%%%%%%%%%%%%%%%%%%%%%%%%%%%%%%%%%%%%%%%%%%

\begin{frame}[fragile] \frametitle{\enquote{Datenstrukturen} in Prolog}
	\footnotesize
	\begin{itemize}
		\item Darstellung von Objekten als Terme über Konstruktoren
		\item keine explizite Deklaration -- implizite Definition über Verwendung in Klauseln
	\end{itemize}
	
	\textbf{natürliche Zahlen:} Prädikat \texttt{nat}
	\vspace{-.75em}
	\begin{itemize} 
		\item nullstelliger Konstruktor \texttt{0} 
		\item einstelliger Konstruktor \texttt{s(X)}
	\end{itemize}
	
	\begin{lstlisting}
nat(0).
nat(s(X)) :- nat(X).
	\end{lstlisting}
	
	\textbf{Listen:} Prädikat \texttt{list}
	\hfill Abkürzung:
	\vspace{-.75em}
	\begin{itemize} 
		\item nullstelliger Konstruktor \texttt{nil} 
		\hfill $\leadsto$ \texttt{[]} 
		\item zweistelliger Konstruktor \texttt{cons(X, Xs)}
		\hfill $\leadsto$ \texttt{[X|Xs]}
	\end{itemize}
	
	\begin{lstlisting}[firstnumber=3]
list(nil).
list(cons(X, Xs)) :- list(Xs).
	\end{lstlisting}
\end{frame}

\begin{frame}[fragile] \frametitle{\enquote{Datenstrukturen} in Prolog}
	\footnotesize
	
	\textbf{Erinnerung:} natürliche Zahlen, Listen
	
	\begin{lstlisting}
nat(0).
nat(s(X)) :- nat(X).
list(nil).
list(cons(X, Xs)) :- list(Xs).
	\end{lstlisting}
	
	\textbf{Bäume:} Prädikat \texttt{istree}
	\vspace{-0.5em}
	\begin{itemize} \footnotesize
		\item nullstelliger Konstruktor \texttt{nil}
		\item dreistelliger Konstruktor \texttt{tree(X,L,R)} 
	\end{itemize}
	
	\begin{lstlisting}[firstnumber=5]
istree(nil).
istree(tree(_, L, R)) :- istree(L), istree(R).
	\end{lstlisting}
\end{frame}




\section{Aufgabe 1 \\ \itshape \normalsize Listen}

\begin{frame}[fragile] \frametitle{Aufgabe 1 -- Teil (a)}
	\footnotesize
	
	\textbf{Ziel:} binäre Relation \texttt{sublist} mit 
	\begin{equation*}
		(\ell_1 , \ell_2) \in \texttt{sublist} 
		\quad \Leftrightarrow \quad
		\ell_1 \subseteq \ell_2
	\end{equation*}
	$\leadsto$ $\ell_1$ soll \textit{Teilliste} von $\ell_2$ sein
	
	\begin{lstlisting}
nat (0).
nat(s(X)) :- nat(X).

listnat ([]).
listnat ([X|XS]) :- nat(X), listnat(XS).

	\end{lstlisting}
	\pause
	\begin{lstlisting}[firstnumber=6]
sublist(Xs   , [Y|Ys]) :- nat(Y), sublist(Xs, Ys). 
sublist(Xs   , Ys )    :- prefix(Xs, Ys). 

prefix([]    , Ys )    :- listnat(Ys). 
prefix([X|Xs], [X|Ys]) :- nat(X), prefix(Xs, Ys).
	\end{lstlisting}
\end{frame}

\begin{frame}[fragile] \frametitle{Aufgabe 1 -- Teil (b)}
	\footnotesize
	\textbf{Belegung 1:}
	
	\begin{lstlisting}[style=refutation, basicstyle=\ttfamily\scriptsize]
          ?-  sublist ([<4>|Xs], [<5>, <4>, <3>]).
          ?-  nat(<5>), sublist ([<4>|Xs], [<4>, <3>]). % 6
          ?-* nat(0), sublist ([<4>|Xs], [<4>, <3>]).   % 2
          ?-  sublist ([<4>|Xs], [<4>, <3>]).           % 1
          ?-  prefix ([<4>|Xs], [<4>, <3>]).            % 7
          ?-  nat(<4>), prefix(Xs , [<3>]).             % 10
          ?-* nat(0), prefix(Xs , [<3>]).               % 2
          ?-  prefix(Xs , [<3>]).                       % 1
{Xs = []} ?-  listnat ([<3>]).                          % 9
          ?-  nat(<3>), listnat ([]).                   % 5
          ?-* nat(0), listnat ([]).                     % 2
          ?-  listnat ([]).                             % 1
          ?-  .                                         % 4
	\end{lstlisting}
	Somit also \texttt{Xs = []}.
\end{frame}

\begin{frame}[fragile] \frametitle{Aufgabe 1 -- Teil (b)}
	\footnotesize
	\textbf{Belegung 2:}
	
	\begin{lstlisting}[style=refutation, basicstyle=\ttfamily\scriptsize]
               ?-  sublist ([<4>|Xs], [<5>, <4>, <3>]).
               ?-  nat(<5>), sublist ([<4>|Xs], [<4>, <3>]). 
                                                        % 6
               ?-* nat(0), sublist ([<4>|Xs], [<4>, <3>]).  
                                                        % 2
               ?-  sublist ([<4>|Xs], [<4>, <3>]).      % 1
               ?-  prefix ([<4>|Xs], [<4>, <3>]).       % 7
               ?-  nat(<4>), prefix(Xs , [<3>]).        % 10
               ?-* nat(0), prefix(Xs , [<3>]).          % 2
               ?-  prefix(Xs , [<3>]).                  % 1
{Xs=[<3>|Xs1]} ?-  nat(<3>), prefix(Xs1 , []).          % 10
               ?-* nat(0), prefix(Xs1 , []).            % 2
               ?-  prefix(Xs1 , []).                    % 1
{Xs1 = []}     ?-  listnat ([]).                        % 9
               ?-  .                                    % 4
	\end{lstlisting}
	
	Somit also \texttt{Xs = [<3>|Xs1] = [<3>]}.
\end{frame}


%%%%%%%%%%%%%%%%%%%%%%%%%%%%%%%%%%%%%%%%%%%%%%%%%%%%%%%%%%%%%%%%%%%%%%%%%%%%%%%%

\section{Aufgabe 2 \\ \itshape \normalsize Bäume}

\begin{frame}[fragile] \frametitle{Aufgabe 2 -- Teil (a)}
	\footnotesize
	Wir wollen einen binären Termbaum auswerten. 
	\begin{lstlisting}
nat (0).
nat(s(X)) :- nat(X).
sum(0, Y, Y) :- nat(Y).
sum(s(X), Y, s(S)) :- sum(X, Y, S).
	\end{lstlisting}

	\textbf{Beispiel:} 
	
	\begin{minipage}{\dimexpr0.35\linewidth-\fboxrule-\fboxsep}		
		\begin{forest}
			[\texttt{minus} [\texttt{<4>}] [\texttt{plus} [\texttt{<2>}] [\texttt{<1>}]]]
		\end{forest}
	\end{minipage}
	\begin{minipage}{\dimexpr0.65\linewidth-\fboxrule-\fboxsep}
		Kodierung über zweistelligen Konstruktoren \texttt{plus} und \texttt{minus} 
		\begin{lstlisting}[style=refutation]
minus( 
	<4>,
	plus(<2>, <1>)
)
		\end{lstlisting}
	\end{minipage}

	$\leadsto \operatorname{minus}(4, \operatorname{plus}(2,1)) = 4 - (2 + 1) = 1$
	
\end{frame}

\begin{frame}[fragile] \frametitle{Aufgabe 2 -- Teil (a)}
	\footnotesize
	
	Wir wollen einen binären Termbaum auswerten. 
	\begin{lstlisting}[basicstyle=\ttfamily\tiny]
nat (0).
nat(s(X)) :- nat(X).
sum(0, Y, Y)       :- nat(Y).
sum(s(X), Y, s(S)) :- sum(X, Y, S).
	\end{lstlisting}
	
	\pause
	
	\begin{lstlisting}[firstnumber=5, basicstyle=\ttfamily\tiny]
eval( X         , X ) :-  nat(X). 
eval( plus (L,R), X ) :-  eval(L, LE), eval(R, RE), sum(LE, RE,  X). 
eval( minus(L,R), X ) :-  eval(L, LE), eval(R, RE), sum(RE,  X, LE).
	\end{lstlisting}
\end{frame}

\begin{frame}[fragile] \frametitle{Aufgabe 2 -- Teil (b)}
	\scriptsize
	\textbf{Gegeben:} zwei Bäume
	\begin{align*}
	\texttt{<t1>} &= \texttt{tree(a, tree(b, nil, nil), tree(v, nil, nil))} \\
	\texttt{<t2>} &= \texttt{tree(c, nil, tree(d, nil, nil))}
	\end{align*}
	\begin{lstlisting}[basicstyle=\ttfamily\tiny]
istree(nil).
istree(tree(_, L, R)) :- istree(L), istree(R).

insert(nil, _, nil).
insert(tree(v, _, _), T, T) :- istree(T).
insert(tree(X, L, R), T, tree(X, LT, RT)) :- insert(L, T, LT), insert(R, T, RT).
	\end{lstlisting}
	\textbf{Anschauung:} \texttt{insert(B, T, X)} ist wahr falls \texttt{X} ensteht, indem alle Knoten \texttt{v} in \texttt{B} durch den Baum \texttt{T} ersetzt werden oder nicht
	
	\textbf{Beispiel:}
	\vspace{-1em}
	\begin{tabular}{ccc}
		\texttt{B} & \texttt{T} & \texttt{X} \\
		\texttt{<t1>} & \texttt{<t2>} & Ergebnis(se) \\
		\scalebox{0.7}{ \adjustbox{valign=t}{
				\begin{forest}{for tree = {l=1em}}
					[\texttt{a} [\texttt{b} [\texttt{nil}] [\texttt{nil}]] [\texttt{v} [\texttt{nil}] [\texttt{nil}]]]
			\end{forest} } }
		&
		\scalebox{0.7}{ \adjustbox{valign=t}{
				\begin{forest}{for tree = {l=1em}}
					[\texttt{c} [\texttt{nil}] [\texttt{d} [\texttt{nil}] [\texttt{nil}]]]
			\end{forest} } }
		&
		\scalebox{0.7}{ \adjustbox{valign=t}{ 
				\begin{forest}{for tree = {l=1em}}
					[\texttt{a} [\texttt{b} [\texttt{nil}] [\texttt{nil}]] [\texttt{v} [\texttt{nil}] [\texttt{nil}]]]
			\end{forest} } 
			\adjustbox{valign=b}{oder}
			\adjustbox{valign=t}{ 
				\begin{forest}{for tree = {l=1em}}
					[\texttt{a} [\texttt{b} [\texttt{nil}] [\texttt{nil}]] [\texttt{c} [\texttt{nil}] [\texttt{d} [\texttt{nil}] [\texttt{nil}]]]]
			\end{forest} } }
	\end{tabular}
\end{frame}

\begin{frame}[fragile] \frametitle{Aufgabe 2 -- Teil (b)}
	\scriptsize
	\textbf{Gegeben:} zwei Bäume
	\begin{align*}
		\texttt{<t1>} &= \texttt{tree(a, tree(b, nil, nil), tree(v, nil, nil))} \\
		\texttt{<t2>} &= \texttt{tree(c, nil, tree(d, nil, nil))}
	\end{align*}
	\begin{lstlisting}[basicstyle=\ttfamily\tiny]
istree(nil).
istree(tree(_, L, R)) :- istree(L), istree(R).

insert(nil, _, nil).
insert(tree(v, _, _), T, T) :- istree(T).
insert(tree(X, L, R), T, tree(X, LT, RT)) :- insert(L, T, LT), insert(R, T, RT).
	\end{lstlisting}
	\textbf{Gesucht:} eine Belegungen für \texttt{X}, die \texttt{?- insert(<t1>, <t2>, X).} erfüllt
	


	\textbf{Alternative 1:}
	\begin{lstlisting}[style=refutation, basicstyle=\ttfamily\tiny]
                           ?-  insert(<t1>, <t2>, X).
{X = tree(a, LT1, RT1)}    ?-  insert(tree(b, nil, nil), <t2>, LT1), 
                               insert(tree(v, nil, nil), <t2>, RT1).   % 6
{LT1 = tree(b, LT2, RT2)}  ?-  insert(nil , <t2>, LT2), 
                               insert(nil, <t2>, RT2), 
                               insert(tree(v, nil, nil), <t2>, RT1).   % 6
{LT2 = nil, RT2 = nil}     ?-* insert(tree(v, nil, nil), <t2>, RT1).   % 4
{RT1 = <t2>}               ?-  istree(<t2>).                           % 5
                           ?-* istree(nil), istree(nil), istree(nil).  % 2
                           ?-* .                                       % 4
	\end{lstlisting}
\end{frame}

\begin{frame}[fragile] \frametitle{Aufgabe 2 -- Teil (b)}
	\scriptsize
	\textbf{Gegeben:} zwei Bäume
	\begin{align*}
	\texttt{<t1>} &= \texttt{tree(a, tree(b, nil, nil), tree(v, nil, nil))} \\
	\texttt{<t2>} &= \texttt{tree(c, nil, tree(d, nil, nil))}
	\end{align*}
	\begin{lstlisting}[basicstyle=\ttfamily\tiny]
	istree(nil).
	istree(tree(_, L, R)) :- istree(L), istree(R).
	
	insert(nil, _, nil).
	insert(tree(v, _, _), T, T) :- istree(T).
	insert(tree(X, L, R), T, tree(X, LT , RT)) 
	:- insert(L, T, LT), insert(R, T, RT).
	\end{lstlisting}
	\textbf{Gesucht:} eine Belegungen für \texttt{X}, die \texttt{?- insert(<t1>, <t2>, X).} erfüllt

	
	\textbf{Alternative 2:} 
	{\tiny die ersten vier Goals stimmen mit Alternative 1 überein}
	\begin{lstlisting}[style=refutation, basicstyle=\ttfamily\tiny\color{cdgray}]
                           ?-  insert(<t1>, <t2>, X).
{X = tree(a, LT1, RT1)}    ?-  insert(tree(b, nil, nil), <t2>, LT1), 
                               insert(tree(v, nil, nil), <t2>, RT1).   % 6
{LT1 = tree(b, LT2, RT2)}  ?-  insert(nil , <t2>, LT2), 
                               insert(nil, <t2>, RT2), 
                               insert(tree(v, nil, nil), <t2>, RT1).   % 6
{LT2 = nil, RT2 = nil}     ?-* insert(tree(v, nil, nil), <t2>, RT1).   % 4
	\end{lstlisting} \vspace{-.5\baselineskip}
	\begin{lstlisting}[style=refutation, basicstyle=\ttfamily\tiny]
{RT1 = tree(v, LT3, RT3)}  ?-  insert(nil, <t2>, LT3),
                               insert(nil, <t2>, RT3).                 % 6
{RT3 = nil, LT3 = nil}     ?-* .                                       % 4
	\end{lstlisting}
	
\end{frame}
%%%%%%%%%%%%%%%%%%%%%%%%%%%%%%%%%%%%%%%%%%%%%%%%%%%%%%%%%%%%%%%%%%%%%%%%%%%%%%%%


\section{Ein weiteres Beispiel \\ \itshape \normalsize aus der Aufgabensammlung}

\setbeamercolor{frametitle}{fg=white, bg=cdgray}	

\begin{frame}[fragile] \frametitle{Aufgabe AGS 13.5 -- Teil (a)}
	\footnotesize
	\textcolor{cdblue}{Gegeben} sei folgender Prolog-Code:
	\begin{lstlisting}
subt( X  , X         ). 
subt( S1 , s(_ , T2) ) :- subt(S1,T2).
subt( S1 , s(T1, _ ) ) :- subt(S1,T1).
	\end{lstlisting}
	\textcolor{cdblue}{Gesucht} sind Belegungen für \texttt{X} und {Y} für das Goal \texttt{?- subt(s(X, Y), s(s(a, b), s(b, a))).}
	\pause
	
	\begin{lstlisting}[style=refutation, basicstyle=\ttfamily\scriptsize]
                        ?-  subt(s(X,Y), s(s(a,b), s(b,a))). 
{X = s(a,b), Y=s(b,a)}  ?-  .                       % 1  

                        ?-  subt(s(X,Y), s(s(a,b), s(b,a))).
                        ?-  subt(s(X,Y), s(b,a)).   % 2
{X = b, Y=a}            ?-  .                       % 1

                        ?-  subt(s(X,Y), s(s(a,b), s(b,a))).
                        ?-  subt(s(X,Y), s(a,b)).   % 3
{X = a, Y=b }           ?-  .                       % 1
	\end{lstlisting}
\end{frame}

\begin{frame}[fragile] \frametitle{Aufgabe AGS 13.5 -- Teil (b)}
	\footnotesize
	\textcolor{cdblue}{Gegeben}  sei folgender Prolog-Code:
	\begin{lstlisting}
subt( X  , X         ). 
subt( S1 , s(_ , T2) ) :- subt(S1,T2).
subt( S1 , s(T1, _ ) ) :- subt(S1,T1).
	\end{lstlisting}
	\textcolor{cdblue}{Gesucht} sind drei Lösungen für das Goal \texttt{?- subt(s(a, a), X).} 
	\pause

	\begin{lstlisting}[style=refutation, basicstyle=\ttfamily\scriptsize]
                  ?-  subt(s(a,a), X). 
{X = s(a,a)}      ?-  .                  % 1
                                         (*@\centering$\Rightarrow$@*) X = s(a,a))
                  ?-  subt(s(a,a), X). 
{X = s(_ , X1)}   ?-  subt(s(a,a), X1).  % 2 
{X1 = s(a,a)}     ?-  .                  % 1 
                                         (*@\centering$\Rightarrow$@*) X = s(a,s(a,a))

                  ?-  subt(s(a,a), X). 
{X = s(X2, _)}    ?-  subt(s(a,a), X2).  % 3 
{X2 = s(a,a) }    ?-  .                  % 1 
                                         (*@\centering$\Rightarrow$@*) X = s(s(a,a),c)
	\end{lstlisting}
\end{frame}

\end{document}

